% Part: first-order-logic
% Chapter: introduction
% Section: soundness-completeness

\documentclass[../../../include/open-logic-section]{subfiles}

\begin{document}

\olfileid[fr]{fol}{int}{scp}

\olsection{Correction et complétude}

Nous présenterons aussi des systèmes de !!{derivation} pour la logique
du premier ordre. Les logiciens en ont élaboré de nombreux, mais tous
définissent la même relation de !!{derivability} entre
!!{sentence}s. Nous disons que $!A$ est \emph{!!{derivable} à partir
de~$\Gamma$}, et nous écrivons $\Gamma \Proves !A$, s'il existe une
!!{derivation} d'un type défini avec précision. Les
!!{derivation}s sont toujours des agencements finis de symboles :
parfois des listes d'!!{sentence}s, parfois des structures plus
complexes. Les systèmes de !!{derivation} doivent fournir un outil qui
permet de déterminer si !!a{sentence} est entraîné par un
ensemble~$\Gamma$. Pour qu'ils remplissent cette fonction, il faut que
$\Gamma \Entails !A$ si et seulement si $\Gamma \Proves !A$.

Si $\Gamma \Proves !A$ mais non $\Gamma \Entails !A$, notre système de
!!{derivation} est trop fort : il démontre trop de choses. La propriété
selon laquelle $\Gamma \Proves !A$ entraîne
$\Gamma \Entails !A$ s'appelle la \emph{correction}; c'est une exigence
minimale pour tout bon système de !!{derivation}. Inversement, si
$\Gamma \Entails !A$ mais non $\Gamma \Proves !A$, notre système de
!!{derivation} est trop faible : il ne démontre pas assez de choses. La
propriété selon laquelle $\Gamma \Entails !A$ entraîne
$\Gamma \Proves !A$ s'appelle la \emph{complétude}. La correction se
démontre généralement assez facilement, par induction sur la structure
des !!{derivation}s, elles-mêmes définies inductivement. La complétude
est plus difficile à établir.

La correction et la complétude ont plusieurs conséquences importantes.
Si l'on peut !!{derive} une contradiction, telle que
$!A \land \lnot !A$, d'un ensemble d'!!{sentence}s~$\Gamma$, celui-ci
est dit \emph{incohérent}. Un ensemble incohérent~$\Gamma$ ne possède
aucun modèle : il est insatisfaisable. La complétude donne la réciproque :
tout ensemble~$\Gamma$ qui n'est pas incohérent, autrement dit qui est
\emph{cohérent}, possède un modèle. Cette affirmation est en fait
équivalente à la complétude, et c'est sous cette forme que nous
démontrerons celle-ci. Le résultat est profond et peut paraître
surprenant : le seul fait qu'on ne puisse pas !!{derive}
$!A \land \lnot !A$ de~$\Gamma$ garantit l'existence
d'!!a{structure} conforme à la description donnée par~$\Gamma$. La
complétude répond donc à la question suivante : quels ensembles
d'!!{sentence}s possèdent des modèles ? Exactement les ensembles
cohérents.

Les théorèmes de correction et de complétude ont deux conséquences
majeures : les théorèmes de compacité et de L\"owenheim--Skolem. Ce sont
des résultats importants de la théorie des modèles, qui servent à en
établir beaucoup d'autres. Nous en avons déjà mentionné deux : la
logique du premier ordre ne peut exprimer ni que le !!{domain}
d'!!a{structure} est fini, ni qu'il est !!{nonenumerable}.

Historiquement, il a fallu beaucoup de temps pour comprendre et mettre
au point tout cet édifice : la définition de la syntaxe et de la
sémantique de la logique du premier ordre, la définition de bons
systèmes de !!{derivation}, la preuve de leur correction et de leur
complétude, et la détermination de ce qui peut ou ne peut pas s'exprimer
dans les langages du premier ordre. Nous savons aujourd'hui comment
procéder, mais l'examen de tous les détails peut encore être déroutant
et laborieux. Il est aussi essentiel, car les méthodes élaborées ici
pour le langage formel de la logique du premier ordre sont employées
partout en logique, en informatique et en linguistique. L'effort
consacré aux détails finit donc par porter ses fruits.

\end{document}
